\begin{abstract}
    \thispagestyle{kapitelkopfzeile}
    \textbf{Abstract:}

    This thesis presents an exploration of the synergies between the worlds of High
    Performance Computing and Cloud Computing, aiming to devise an interactive
    exploratory computing environment within a cloud-native setting capable of
    handling HPC-scale data analysis tasks.
    The research addresses the challenge of merging the computational power and
    efficiency of HPC systems with the versatility and flexibility of Cloud Computing
    and the need for an  easy-to-use and user-friendly interface;
    which enables real-time research in a familiar programming environment.
    For this purpose, an architecture, was designed from literature by way of
    a Conceptual Deductive Analysis, and the resulting system design was implemented
    in an iterative development process applying the gained insights to update
    the system design with the experimental knowledge.
    The artifact created in this way is a bare metal Kubernetes cluster, running
    a fleet of Arkouda worker instances, which can be interfaced with via a
    Jupyter Notebook with native a native integration providing a seamless
    experience for the user.
    The goal for this artifact being to extend the capabilities of the Pachykuda
    system designed in the previous work by the author, extending the
    GitOps based pipeline workflow orchestrator with the ability interactively
    explore and build data analysis pipelines within a Jupyter environment.
    This thesis contributes to the field by constructing an experimentally
    verified design for the integration of a computational notebook interfaces
    with Kubernetes and highly parallel compute frameworks and by providing a tangible artifact.


\end{abstract}